\documentclass[aspectratio=169,12pt]{beamer}

% Compile with XeLaTeX, LuaLaTeX, or Tectonic.
% A wrapper may redefine \DeckFontSetup before inputting this file.
\usepackage[UTF8,scheme=plain,fontset=none]{ctex}
\usepackage{booktabs,array}

\providecommand{\DeckFontSetup}{%
  \setsansfont[
    Path={fonts/noto/}, Extension=.otf,
    UprightFont=*-Regular, BoldFont=*-Bold
  ]{NotoSansSC}%
  \setmainfont[
    Path={fonts/noto/}, Extension=.otf,
    UprightFont=*-Regular, BoldFont=*-Bold
  ]{NotoSansSC}%
  \setCJKsansfont[
    Path={fonts/noto/}, Extension=.otf,
    UprightFont=*-Regular, BoldFont=*-Bold
  ]{NotoSansSC}%
  \setCJKmainfont[
    Path={fonts/noto/}, Extension=.otf,
    UprightFont=*-Regular, BoldFont=*-Bold
  ]{NotoSansSC}%
}
\DeckFontSetup

\usetheme{Madrid}
\usefonttheme{professionalfonts}

\definecolor{BeamerRoyal}{HTML}{3939B8}
\definecolor{BeamerDeep}{HTML}{171B72}
\definecolor{BeamerInk}{HTML}{20212B}
\definecolor{SoftBlue}{HTML}{F1F2FF}
\definecolor{SoftGray}{HTML}{F5F5F7}
\definecolor{WarmAccent}{HTML}{B2472E}

\setbeamercolor{normal text}{fg=BeamerInk,bg=white}
\setbeamercolor{structure}{fg=BeamerRoyal}
\setbeamercolor{frametitle}{fg=white,bg=BeamerRoyal}
\setbeamercolor{title}{fg=white,bg=BeamerRoyal}
\setbeamercolor{subtitle}{fg=BeamerInk}
\setbeamercolor{block title}{fg=white,bg=BeamerRoyal}
\setbeamercolor{block body}{fg=BeamerInk,bg=SoftBlue}
\setbeamercolor{alertblock title}{fg=white,bg=WarmAccent}
\setbeamercolor{alertblock body}{fg=BeamerInk,bg=WarmAccent!9}
\setbeamercolor{author in head/foot}{fg=white,bg=BeamerDeep}
\setbeamercolor{title in head/foot}{fg=white,bg=BeamerRoyal}
\setbeamercolor{date in head/foot}{fg=white,bg=BeamerDeep}

\setbeamertemplate{items}[ball]
\setbeamertemplate{blocks}[rounded][shadow=true]
\setbeamertemplate{navigation symbols}{}
\setbeamertemplate{headline}{}

\setbeamertemplate{title page}{%
  \vspace{0.90cm}
  \begin{beamercolorbox}[wd=.80\paperwidth,sep=.24cm,center,rounded=true,shadow=true]{title}
    \usebeamerfont{title}\inserttitle
  \end{beamercolorbox}
  \vspace{0.72cm}
  \begin{center}
    {\usebeamerfont{subtitle}\insertsubtitle\par}
    \vspace{0.86cm}
    {\usebeamerfont{author}\insertauthor\par}
    \vspace{0.24cm}
    {\usebeamerfont{institute}\insertinstitute\par}
    \vspace{0.74cm}
    {\usebeamerfont{date}\insertdate\par}
  \end{center}
}

\setbeamertemplate{frametitle}{%
  \nointerlineskip
  \begin{beamercolorbox}[wd=\paperwidth,ht=.64cm,dp=.14cm,leftskip=.38cm]{frametitle}
    \usebeamerfont{frametitle}\insertframetitle
  \end{beamercolorbox}
}

\setbeamertemplate{footline}{%
  \leavevmode
  \hbox{%
    \begin{beamercolorbox}[wd=.31\paperwidth,ht=2.6ex,dp=1.05ex,center]{author in head/foot}
      \usebeamerfont{author in head/foot}\insertshortauthor\enspace(\insertshortinstitute)
    \end{beamercolorbox}%
    \begin{beamercolorbox}[wd=.36\paperwidth,ht=2.6ex,dp=1.05ex,center]{title in head/foot}
      \usebeamerfont{title in head/foot}\insertshorttitle
    \end{beamercolorbox}%
    \begin{beamercolorbox}[wd=.33\paperwidth,ht=2.6ex,dp=1.05ex,center]{date in head/foot}
      \usebeamerfont{date in head/foot}\insertshortdate\hspace{1.2em}
      \insertframenumber/\inserttotalframenumber
    \end{beamercolorbox}%
  }
}

\setbeamersize{text margin left=1.12cm,text margin right=1.12cm}
\setbeamerfont{title}{size*={21pt}{26pt},series=\bfseries}
\setbeamerfont{subtitle}{size*={13.5pt}{18pt}}
\setbeamerfont{author}{size*={12.5pt}{16.5pt},series=\bfseries}
\setbeamerfont{institute}{size*={10.5pt}{14pt}}
\setbeamerfont{date}{size*={11.5pt}{15pt}}
\setbeamerfont{frametitle}{size*={17pt}{20.5pt}}
\setbeamerfont{normal text}{size*={11.8pt}{16pt}}
\setbeamerfont{itemize/enumerate body}{size*={11.6pt}{16pt}}
\setbeamerfont{itemize/enumerate subbody}{size*={10.7pt}{14.5pt}}
\setbeamerfont{block title}{size*={11.5pt}{14.8pt},series=\bfseries}
\setbeamerfont{block body}{size*={11.1pt}{15pt}}
\setbeamerfont{author in head/foot}{size*={7pt}{8.5pt}}
\setbeamerfont{title in head/foot}{size*={7pt}{8.5pt}}
\setbeamerfont{date in head/foot}{size*={7pt}{8.5pt}}

\newcommand{\key}[1]{\textcolor{BeamerRoyal}{\textbf{#1}}}
\newcommand{\takeaway}[1]{%
  \begin{beamercolorbox}[wd=\linewidth,sep=.16cm,center,rounded=true,shadow=true]{block body}
    \textcolor{BeamerDeep}{\textbf{#1}}
  \end{beamercolorbox}
}
\newcommand{\bigarrow}{\textcolor{BeamerRoyal}{\large\bfseries →}}

\title[NV-Segment-CTMR]{NV-Segment-CTMR 模型简介}
\subtitle{从 CT / MRI 到三维器官标签}
\author[学习笔记]{NV-Segment-CTMR 学习笔记}
\institute[医学影像分割]{面向入门学习者的简明讲解}
\date[2026.09]{2026 年 9 月}

\begin{document}

\begin{frame}
  \titlepage
\end{frame}

\begin{frame}{它解决什么问题？}
  \vspace{0.12cm}
  \begin{columns}[c,onlytextwidth]
    \column{.28\textwidth}
      \begin{block}{输入}
        \centering
        三维 CT / MRI\\[0.08cm]
        {\small 每个 voxel 是图像强度}
      \end{block}
    \column{.06\textwidth}\centering\bigarrow
    \column{.28\textwidth}
      \begin{block}{模型}
        \centering
        自动寻找结构\\[0.08cm]
        {\small 器官、组织等类别}
      \end{block}
    \column{.06\textwidth}\centering\bigarrow
    \column{.28\textwidth}
      \begin{block}{输出}
        \centering
        三维标签体数据\\[0.08cm]
        {\small 每个 voxel 一个标签}
      \end{block}
  \end{columns}

  \vspace{0.30cm}
  \begin{itemize}
    \item 支持 CT 与 MRI（身体和脑部），用于\key{自动分割}
    \item 官方类别表包含 \key{345+ 类}；不同模态使用的类别集合并不完全相同
    \item 没有属于目标结构的位置会保留为 \textbf{background}
  \end{itemize}
  \vfill
  \takeaway{输入是图像强度，输出才是器官标签。}
\end{frame}

\begin{frame}{一张图看懂整体流程}
  \vspace{0.10cm}
  \begin{columns}[T,onlytextwidth]
    \column{.31\textwidth}
      \begin{block}{1. Preprocessing}
        统一方向、spacing 与强度，让不同扫描更容易比较。
      \end{block}
    \column{.31\textwidth}
      \begin{block}{2. Sliding Window}
        把很大的三维影像分成 GPU 能处理的小块。
      \end{block}
    \column{.31\textwidth}
      \begin{block}{3. 3D CNN}
        编码图像，再逐步恢复空间位置与细节。
      \end{block}
  \end{columns}

  \vspace{0.05cm}
  \begin{columns}[T,onlytextwidth]
    \column{.31\textwidth}
      \begin{block}{4. Feature Map}
        把灰度变成描述边缘、形状和位置的多维特征。
      \end{block}
    \column{.31\textwidth}
      \begin{block}{5. Class Matching}
        图像特征与器官的 class embedding 进行匹配。
      \end{block}
    \column{.31\textwidth}
      \begin{block}{6. Merge \& Output}
        融合重叠窗口，恢复坐标并生成最终标签。
      \end{block}
  \end{columns}

  \vfill
  \takeaway{{\fontsize{11.8}{15.2}\selectfont
    先整理数据，再提取特征、寻找类别，最后拼回完整结果。}}
\end{frame}

\begin{frame}{Preprocessing：先把数据整理好}
  \begin{columns}[T,onlytextwidth]
    \column{.47\textwidth}
      \begin{block}{Orientation｜统一方向}
        把影像调整到一致的坐标方向，避免左右、前后定义混乱。
      \end{block}
      \begin{block}{Spacing｜统一物理尺度}
        重采样到统一的 voxel 尺度，让相同器官看起来大小更接近。
      \end{block}
    \column{.47\textwidth}
      \begin{block}{Crop｜裁掉大块背景}
        保留包含身体的区域，减少没有信息的计算。
      \end{block}
      \begin{block}{Intensity｜处理强度}
        缩放并截断极端值，让不同扫描的数值范围更稳定。
      \end{block}
  \end{columns}

  \vspace{0.10cm}
  \takeaway{预处理只是在“整理试卷”，真正的器官识别还没有开始。}
\end{frame}

\begin{frame}{训练：模型如何一步步变准？}
  \begin{columns}[T,onlytextwidth]
    \column{.61\textwidth}
      \begin{enumerate}
        \item \key{前向传播}：输入影像，模型按当前参数生成预测 mask
        \item \key{比较}：预测与人工标注的 ground-truth mask 产生 loss
        \item \key{反向传播}：计算各参数对错误应负多少责任
        \item \key{更新}：optimizer 根据梯度修改参数，再重复训练
      \end{enumerate}
    \column{.34\textwidth}
      \begin{block}{可以记成五个词}
        \centering
        尝试 → 比较\\[0.10cm]
        {\color{BeamerRoyal}\Large ↓}\\[-0.02cm]
        修改 ← 追责\\[0.10cm]
        {\color{BeamerRoyal}\Large ↓}\\[-0.02cm]
        再尝试
      \end{block}
  \end{columns}

  \vfill
  \takeaway{{\fontsize{11.8}{15.2}\selectfont
    训练就是不断“先做题、看答案、找原因、改方法”。}}
\end{frame}

\begin{frame}{Sliding Window：为什么要切成小块？}
  \begin{columns}[T,onlytextwidth]
    \column{.47\textwidth}
      \begin{block}{主要原因：显存有限}
        完整三维 CT / MRI 往往很大，难以一次全部送入 GPU。滑窗让模型一次只处理一个 3D patch。
      \end{block}
    \column{.47\textwidth}
      \begin{block}{为什么窗口要重叠？}
        同一区域会被多个窗口预测，再把结果融合，可减轻 patch 边缘的不稳定与断裂。
      \end{block}
  \end{columns}

  \vspace{0.22cm}
  \begin{alertblock}{最容易讲错的一点}
    滑动窗口不是为了让 patch 直接“互相交流”，而是为了控制单次显存；重叠主要用于减少块边缘伪影。
  \end{alertblock}

  \begin{itemize}
    \item 整幅图基本仍要被扫过，重叠区域还会重复计算
    \item 所以它主要降低\key{单次内存需求}，不一定降低总计算量
  \end{itemize}
\end{frame}

\begin{frame}{3D Encoder--Decoder 与 Feature Map}
  \begin{columns}[c,onlytextwidth]
    \column{.27\textwidth}
      \begin{block}{Encoder}
        \centering
        缩小空间范围\\
        看更大的上下文
      \end{block}
    \column{.07\textwidth}\centering\bigarrow
    \column{.31\textwidth}
      \begin{block}{Feature Map}
        \centering
        每个位置不再只有灰度，\\
        而是一组图像特征
      \end{block}
    \column{.07\textwidth}\centering\bigarrow
    \column{.27\textwidth}
      \begin{block}{Decoder}
        \centering
        恢复空间分辨率\\
        找回边界与位置
      \end{block}
  \end{columns}

  \vspace{0.18cm}
  \begin{itemize}
    \item 这些特征可以描述边缘、纹理、形状、位置和周围结构
    \item 本模型采用 SegResNet 式 3D 编码器--解码器；中间特征图会缩小，再由 decoder 恢复空间位置
  \end{itemize}
  \vfill
  \takeaway{{\fontsize{11.8}{15.2}\selectfont
    Feature Map 是模型整理出的“线索地图”，还不是最终答案。}}
\end{frame}

\begin{frame}{Class Embedding：告诉模型要找什么}
  \begin{columns}[T,onlytextwidth]
    \column{.45\textwidth}
      \begin{block}{Image Feature}
        \centering
        “这个位置看起来有什么？”\\[0.10cm]
        {\small 来自 3D 图像的特征}
      \end{block}
    \column{.45\textwidth}
      \begin{block}{Class Embedding}
        \centering
        “现在要寻找哪一类？”\\[0.10cm]
        {\small 每一类都有一个可学习编码}
      \end{block}
  \end{columns}

  \vspace{0.14cm}
  \begin{center}
    \bigarrow\quad
    \textbf{匹配后得到该类别的分数与 mask}
  \end{center}

  \begin{itemize}
    \item 模型通过类别 ID 选择相应的 class embedding
    \item 图像特征与类别编码匹配，为所选模态的支持类别生成 mask
    \item 它学到的是\key{类别编码与图像特征的对应关系}，不是词语的字面含义
  \end{itemize}
\end{frame}

\begin{frame}{最终输出：三维 Label Volume}
  \begin{columns}[c,onlytextwidth]
    \column{.20\textwidth}
      \begin{block}{Patch 预测}
        \centering 各窗口的器官 mask
      \end{block}
    \column{.045\textwidth}\centering\bigarrow
    \column{.20\textwidth}
      \begin{block}{融合}
        \centering 合并重叠位置的预测
      \end{block}
    \column{.045\textwidth}\centering\bigarrow
    \column{.20\textwidth}
      \begin{block}{恢复}
        \centering 映射回原影像坐标
      \end{block}
    \column{.045\textwidth}\centering\bigarrow
    \column{.22\textwidth}
      \begin{block}{Label Volume}
        \centering 每个 voxel 一个类别标签
      \end{block}
  \end{columns}

  \vspace{0.20cm}
  \begin{itemize}
    \item 最终输出不是 feature vector，而是器官 / 组织标签或 background
    \item 标签体数据可用于着色、统计体积，并进一步生成三维表面模型
  \end{itemize}

  \vspace{0.08cm}
  \takeaway{{\fontsize{11.3}{14.8}\selectfont
    滑窗解决“图像太大”；Feature Map 说明“看到了什么”；\\
    Class Embedding 说明“要找什么”；Label Volume 给出最终答案。}}
\end{frame}

\end{document}
